\subsection{Frontera exacta y escalar reducido}\label{sec:B4T04}
Se usan las Secciones~\ref{sec:B4T01}--\ref{sec:B4T03}, con $z=\bar p-w$, $u=d-D^Tz$ y
\[
 \mathcal K(a,t)=(a+t)\log\frac{a+t}{a}-t,
 \qquad \mathcal K(a,-a)=a.
\]
Las convenciones de $0\log0$ se toman por continuidad.
El padre es estacionario para $H_B$, no para $H_B+\ell$.

\paragraph{Identidad de acción sin desarrollo de Taylor en las casillas eliminadas.}
Extendamos $\bar p$ al soporte padre poniendo uno en $c$ y cero en
$F\setminus\{c\}$. El incremento conserva todas las restricciones RC,
de modo que el término lineal del gradiente parental se anula exactamente.
La identidad de Bregman de cada sumando de $H_B$ da
\[
 \boxed{\begin{aligned}
 H_{\rm hijo}(\bar p)-H_{\rm padre}(p)-\log p_c
 ={}&-3\sum_{F\setminus c}(\log q_e+p_e)
 -\sum_{E'}\mathcal K(p_e,z_e)\\
 &+3\sum_{E'}\mathcal K(q_e,-z_e)
 -\sum_{D\cap I}\mathcal K(h_d,u_d).
 \end{aligned}}
\]
Los términos de frontera se verifican directamente: la casilla $c$ aporta $\log p_c+4q_c$,
las demás casillas eliminadas aportan $-3\log q_e-4p_e$ y las dos
diagonales eliminadas aportan $-h_{D_+}-h_{D_-}$.
Se usa la identidad exacta
$\sum_{F\setminus c}p_e=4q_c-h_{D_+}-h_{D_-}$.
Las cuatro líneas ordinarias sólo se cortan en $c$.

\paragraph{Cancelación local de primer orden.}
Con $b(p)=(1-4p)/(1-p)^4$, el factor local y el bloque retirado dan
\[
 -3\sum_{F\setminus c}(\log q_e+p_e)
 +\tfrac12\sum_F\log b(p_e)+\tfrac12\log\det S_{JJ}=O(k^{-2}).
\]
En efecto, los tres términos principales son, respectivamente,
$\tfrac32\sum_Fp_e^2$, $-3\sum_Fp_e^2$ y $\tfrac32\sum_Fp_e^2$.
La multiplicidad cuatro de $c$ sólo produce un cambio $O(p_c^2)=O(k^{-2})$.
Los restos se acotan por $\sum_Fp_e^3=O(k^{-2})$ y
$\|S_{JJ}-I\|^2=O(k^{-2})$, pues $J$ tiene tamaño cuatro.

\paragraph{Expansión en el punto reequilibrado.}
Sean $L=H+\ell+\tfrac12\log(2\pi)$ en ambos soportes conexos y
\[
 R_{c0}=L_{\rm hijo}(\bar p)-L_{\rm padre}(p)-\log p_c.
\]
Un desarrollo de Taylor sólo en las coordenadas supervivientes y en las vacancias positivas,
junto con la Sección~\ref{sec:B4T01}, produce
\[
 \boxed{R_{c0}=-\tfrac12\left(\sum_{E'}\frac{z_e^2}{p_e}
                  +\sum_{D\cap I}\frac{u_d^2}{h_d}\right)
 +6\sum_{E'}p_ez_e-\tfrac12\operatorname{tr}W_c
 +\tfrac14\operatorname{tr}W_c^2+O(k^{-2}).}
\]
Se han usado $\sum z^2=O(k^{-2})$, $\sum z^2/p=O(k^{-1})$ y
$\sum u^2/h=O(k^{-1})$. Los restos cúbicos son $O(k^{-2})$.
El cambio del factor local en las casillas supervivientes satisface
$-\tfrac12\sum_{E'}\log[b(p+z)/b(p)]=6\sum pz+O(k^{-2})$.

También se tiene $\operatorname{tr}W_c=O(k^{-1})$:
la diagonal de $\mathcal E$ es $O(k^{-2})$ después de corregir las masas RC,
$\|\mathcal E\|_\infty=O(k^{-1})$, y
$P_c^+=\Pi_I^\perp G_{II}\Pi_I^\perp=I+O_{\max}(k^{-1})$.
Aquí $\Pi_I$ proyecta sobre $v_I$; la igualdad se verifica resolviendo
$S(a,y)=(0,f)$ para $f\perp v_I$ y proyectando $y$.
Hay $O(k)$ líneas, por lo que $\sum_{ij}|\mathcal E_{ij}|=O(1)$.

\paragraph{Escalar reducido.}
Sean $e$ el residuo de la Sección~\ref{sec:B4T02}, $\bar R$ su respuesta, $\bar g=\nabla\ell(\bar p)$
y $\delta_c=(R_{\rm padre}\nabla\ell(p))_c/p_c$.
Las Secciones~\ref{sec:B4T02}--\ref{sec:B4T03} y $\log(1+\delta_c)=\delta_c+O(k^{-2})$ dan
\begin{equation}\label{eq:rc}
 r_c=L_{\rm hijo}(p^*)-L_{\rm padre}(p)-\log[p_c(1+\delta_c)]
       =\mathcal C_c+O(k^{-2}),
\end{equation}
\begin{equation}\label{eq:Cc}
 \boxed{\begin{aligned}
 \mathcal C_c={}&-\tfrac12\left(\sum_{E'}z_e^2/p_e+
                                \sum_{D\cap I}u_d^2/h_d\right)
 +6\sum_{E'}p_ez_e\\
 &-\tfrac12\operatorname{tr}W_c+\tfrac14\operatorname{tr}W_c^2
 +\tfrac12e^T\bar R e+\bar g^T\bar R e-\delta_c.
 \end{aligned}}
\end{equation}
Todos sus términos se definen a partir del padre, del reequilibrado de la Sección~\ref{sec:B3T07} y
de un paso de Newton, sin resolver de nuevo el hijo. Estas estimaciones dan
$\mathcal C_c=O(k^{-1})$. Las Secciones~\ref{sec:B4T06}--\ref{sec:B4T09} prueban
$\mathcal C_c=O(k^{-2})$ mediante la resolvente de diagonal uno.
Esta cancelación total es distinta de la cancelación local probada más arriba.

Reunimos lo probado en esta sección.
\begin{lemasub}[Escalar reducido]\label{lem:escalar}
En las hipótesis de la Sección~\ref{sec:B4T01}, para $k$ suficientemente grande vale~\eqref{eq:rc},
con $\delta_c$ como arriba y $\mathcal C_c$ dado por~\eqref{eq:Cc}; además $\mathcal C_c=O(k^{-1})$.
\end{lemasub}
